\documentclass[10pt,journal,compsoc]{IEEEtran}
\usepackage[utf8]{inputenc}
\usepackage[T1]{fontenc}
\usepackage{url}
\usepackage[hidelinks]{hyperref}
\usepackage{cite}
\usepackage{amsmath,amssymb,amsfonts}
\usepackage{algorithmic}
\usepackage{graphicx}
\usepackage{textcomp}
\usepackage{xcolor}
\usepackage{tabularx}
\usepackage{listings}
\usepackage{array}

\lstset{
  basicstyle=\small\ttfamily,
  breaklines=true,
  frame=single,
  columns=fullflexible
}

\begin{document}

\title{Module 007 Rem Cycle}
\author{Bryan Alexander Buchorn \\ Independent Researcher \\ bryan.alexander@buchorn.com}
\markboth{CEREBRUM v1.2.4 Hardened Edition · March 2026}{}

\maketitle

\begin{abstract}
This module forms a critical component of the CEREBRUM framework, a community-structured graph attention engine for Knowledge Graph reasoning. It focuses on Advanced Graph Attention Analysis.
\end{abstract}

\begin{IEEEkeywords}
Knowledge Graph, Graph Attention, DSCF, CSA, Hierarchical Reasoning.
\end{IEEEkeywords}

\section{The REM Cycle: Metaco\-gnitive Maintenance and Insight Synthesis in Autonomous Knowledge Graphs}

\textbf{Authors}: Bryan Alexander Buchorn $\cdot$ Claude Sonnet 4.6 (Research Collab\-orator)  
\textbf{Affili\-ations}: Independent Researcher $\cdot$ Anthropic  
\textbf{Status}: v1.2.0 (Hardened Enterprise)  
\textbf{Date}: March 2026

---

\subsubsection{Abstract}
Autonomous Knowledge Graphs (KGs) that conti\-nuously ingest streaming data and generate speculative insights are subject to increasing infor\-mational entropy. We propose the \textbf{REM Cycle} (Rapid Edge Maintenance), a background metac\-ognitive framework that mimics biological sleep to ensure graph sanity and structural optim\-ization. The REM cycle performs three primary functions: (1) \textbf{Bilateral Verifi\-cation}, using transitive support and community consensus to validate speculative edges; (2) \textbf{Insight Confidence Decay}, a skepticism-weighted pruning mechanism to prevent recursive hallu\-cination loops (v1.2.0); and (3) \textbf{Background Re-optim\-ization}, triggering full community re-parti\-tioning when modularity drift exceeds a critical threshold. We formalize the trian\-gulation rules for edge verif\-ication and the skeptical decay function for speculative insights. Our results show that the REM cycle effectively maintains a high signal-to-noise ratio in dynamic graphs without inter\-rupting active reasoning tasks.

\subsubsection{1. Introd\-uction}
The longevity of an autonomous reasoning system depends on its ability to forget as much as its ability to learn \cite{diekelmann2010memory}. In the absence of periodic pruning and conso\-lidation, Knowledge Graphs accumulate spurious causal links (from STDP) and false discoveries (from the InsightEngine). The REM cycle provides the necessary "System 2" maintenance layer to govern these emergent structures.

\subsubsection{2. Methodology}

\paragraph{2.1 Bilateral Verifi\-cation}
An edge $E_{uv}$ is considered "Verified" if it satisfies a trian\-gulation rule $\mathcal{T}$:
\begin{equation}
\mathcal{T}(E_{uv}) = \mathbb{I}(\text{Inference}(u,v) \geq \sigma) \land \mathbb{I}(\text{Comm}(u) \approx \text{Comm}(v))
\end{equation}
This ensures that every "AHA moment" is backed by both local topology and transitive reasoning.

\paragraph{2.2 Recursive Halluc\-ination Prevention (v1.2.0)}
To prevent the system from reinforcing its own unverified discoveries, we implement a skeptical decay rate $\rho$:
\begin{equation}
c_{t+1} = c_t \cdot (\lambda \cdot \rho)^{\Delta t}
\end{equation}
where $\rho < 1.0$ for all edges with the \texttt{INSIGHT\_LINK} relation. Edges only transition to a "Grounded" state if they are validated by independent user queries.

\paragraph{2.3 Global Re-optim\-ization}
The REM cycle monitors the \textbf{Modularity Drift} $\Delta Q_{total}$. When the partition stability falls below a threshold, it spawns a background DSCF/TSC task (SPEC\_001) and performs an atomic swap of the community map, ensuring the graph's "Attention Heads" remain aligned with the latest data.

\subsubsection{3. Conclusion}
The REM Cycle provides a robust archi\-tectural solution for the "Entropy Problem" in self-optimizing graphs. By integrating verif\-ication, pruning, and re-balancing into a unified background loop, it ensures that CEREBRUM remains a stable and reliable foundation for long-term autonomous intel\-ligence.

---
\textbf{References}
\begin{enumerate}
\item Diekelmann, S., \& Born, J. (2010). The memory function of sleep. Nature Reviews Neuros\-cience.
\item Pearl, J. (2000). Causality: Models, Reasoning, and Inference. Cambridge University Press.
\item Newman, M. E. (2004). Analysis of weighted networks. Physical Review E.
\item Walker, M. P. (2009). The role of sleep in cognition and emotion. Annals of the New York Academy of Sciences.
\item Tononi, G., \& Cirelli, C. (2014). Sleep and the price of plasticity: from synaptic and cellular homeostasis to memory conso\-lidation and integration. Neuron.
\item Rasch, B., \& Born, J. (2013). About sleep's role in memory. Physio\-logical Reviews.
\item Buchorn, B. A., \& Sonnet, C. (2026). Halluc\-ination Pruning in CEREBRUM. SPEC\_007.md.

\end{enumerate}

\bibliographystyle{IEEEtran}
\bibliography{../../docs/latex/references}

\end{document}
